\capituloPregunta{¿Cómo explorar varios factores sin hacer todos los experimentos?}

\section{Motivación}

Cuando existen muchos factores, un diseño factorial completo puede requerir demasiadas corridas. Los diseños Taguchi permiten explorar sistemas con menos tratamientos, siempre que se interpreten con prudencia.

\section{Caso integrado: violeta africana}

El material del curso contiene una práctica con hojas de violeta africana en medios gelificados. Los factores fueron café soluble, miel y azúcar. Se utilizó una matriz Taguchi \(L_4(2^3)\).

La intención inicial era observar propagación vegetal. Sin embargo, al día 13 no aparecieron raíces ni hojas nuevas. En cambio, surgieron respuestas visuales relevantes:

\begin{itemize}
  \item proliferación fúngica;
  \item halo café;
  \item condensación;
  \item oxidación;
  \item pérdida de tejido verde;
  \item necrosis.
\end{itemize}

\begin{idea}
Un experimento también enseña cuando el resultado esperado no ocurre. La clave es transformar el hallazgo inesperado en una nueva pregunta científica.
\end{idea}

\section{Matriz reducida}

Un diseño completo con tres factores a dos niveles requeriría:

\[
2^3 = 8
\]

tratamientos. La matriz Taguchi \(L_4\) reduce el número de combinaciones, conservando una estructura ordenada de exploración.

\section{Caso complementario: datos continuos a niveles}

Otro material del curso muestra cómo convertir variables continuas en niveles discretos mediante cuantiles para reconstruir un diseño tipo Taguchi. Esta idea permite discutir un punto fundamental: no todos los datos nacen como diseño experimental, pero algunos pueden organizarse para exploración si se declaran sus límites.

\section{Errores frecuentes}

\begin{cuidado}
Un Taguchi exploratorio no debe venderse como conclusión definitiva. Reduce corridas y ayuda a detectar patrones, pero sus resultados deben interpretarse según la calidad de medición, el balance y el objetivo del estudio.
\end{cuidado}

\section{Actividad}

Proponga tres factores con dos niveles para estudiar un sistema biológico o de materiales. Construya una matriz reducida de cuatro tratamientos y defina al menos dos variables respuesta.

\section{Conexión}

Cuando la respuesta es continua y se desea estimar tendencias, el siguiente paso es modelar.
